%=============================================================================
% SECTION 8C EQUATION EXTRACTION AND COMPUTATIONAL CHAIN ANALYSIS
% Source: section_alpha_locked.tex + appendix_K.tex (supporting derivation)
% Purpose: Map every equation, identify the computational chain to 137.03599985
% Date: 2026-03-22
%=============================================================================

\section*{Section 8C: Complete Equation Catalog}

%-----------------------------------------------------------------------------
\subsection*{A. Every Equation in section\_alpha\_locked.tex}
%-----------------------------------------------------------------------------

% NOTE: Section 8C does NOT use 8C.1, 8C.2, ... numbering.
% Equations use descriptive labels. We assign sequential catalog numbers below.

\paragraph{Eq.\ [1] --- Connection Laplacian (eq:connection-laplacian, line 26)}
\[
  P = -g^{ij}\nabla_i\nabla_j
\]
\begin{itemize}
  \item \textbf{Computes:} The operator acting on sections of the internal bundle
        over $X = \mathbb{C}P^2 \times S^3$.
  \item \textbf{Input:} Internal metric $g^{ij}$, connection $\nabla$.
  \item \textbf{Output:} Laplace-type operator $P$.
  \item \textbf{Analytic/Numerical:} Analytic (operator definition).
  \item \textbf{Value:} N/A (defines the operator; no numerical value).
\end{itemize}

\paragraph{Eq.\ [2] --- Spectral Action (eq:spectral-action-plateau, line 42)}
\[
  S = \mathrm{Tr}\, f(P/\Lambda^2)
\]
\begin{itemize}
  \item \textbf{Computes:} The spectral action with plateau cutoff $f$.
  \item \textbf{Input:} Operator $P$, cutoff scale $\Lambda$, function $f$.
  \item \textbf{Output:} Action functional $S$.
  \item \textbf{Analytic/Numerical:} Analytic (formal trace).
  \item \textbf{Value:} N/A (functional, not a number).
\end{itemize}

\paragraph{Eq.\ [3] --- Plateau Condition (line 49)}
\[
  f_{-2} = f_{-4} = \cdots = 0
\]
\begin{itemize}
  \item \textbf{Computes:} Constraint on cutoff function moments.
  \item \textbf{Purpose:} Eliminates $a_6, a_8, \ldots$ heat-kernel contributions
        as hidden tuning knobs.
  \item \textbf{Analytic/Numerical:} Analytic (constraint).
\end{itemize}

\paragraph{Eq.\ [4] --- Hilbert Space Dimension (eq:d-dimension, line 66)}
\[
  d = \dim H^0(\mathbb{C}P^1, \mathcal{O}(m)) = m + 1 = k + 4
\]
\begin{itemize}
  \item \textbf{Computes:} Dimension of the finite Hilbert space from Toeplitz
        quantization at level $m = k+3$.
  \item \textbf{Input:} $k = k_{\max} = 60$, shift $+3$ from Spin$^c$ structure.
  \item \textbf{Output:} $d = 64$.
  \item \textbf{Analytic/Numerical:} Analytic (algebraic geometry).
  \item \textbf{Value:} $d = 64$.
\end{itemize}

\paragraph{Eq.\ [5] --- Canonical Bundle (line 72)}
\[
  K_{\mathbb{C}P^2} = \mathcal{O}(-3) \quad\Rightarrow\quad
  L_{\det} = K^{-1} = \mathcal{O}(3)
\]
\begin{itemize}
  \item \textbf{Computes:} The origin of the $+3$ shift ($m = k+3$) from the
        Spin$^c$ determinant line.
  \item \textbf{Analytic/Numerical:} Analytic (rigorous algebraic geometry).
  \item \textbf{Value:} Shift $= 3$.
\end{itemize}

\paragraph{Eq.\ [6] --- Spectral Cutoff / Finite-$k$ Factor (eq:lambda-cutoff, line 79)}
\[
  \Lambda^3 = k \cdot \frac{d-1}{d} = k \cdot \frac{k+3}{k+4}
\]
\begin{itemize}
  \item \textbf{Computes:} The spectral cutoff from determinant-channel removal
        at finite $d$.
  \item \textbf{Input:} $k = 60$, $d = 64$.
  \item \textbf{Output:} $\Lambda^3 = 60 \times (63/64) = 59.0625$.
  \item \textbf{Analytic/Numerical:} Analytic.
  \item \textbf{Value:} $\Lambda^3 = 59.0625$.
  \item \textbf{NOTE:} The text at line 210 also quotes
        $\Lambda^3 = 885.9375$ from ``$k = 60$, $a = 9$, $n = 5$, $N = 3$''.
        This is a DIFFERENT formula:
        $\Lambda^3 = k \cdot (k+3)/(k+4) \times N_{\rm gen} \times n \times \ldots$
        or possibly $k \cdot (k+3)/(k+4) \times 15 = 59.0625 \times 15 = 885.9375$.
        The exact formula producing 885.9375 is NOT displayed as a numbered
        equation in this section. It is stated only in a paragraph.
\end{itemize}

\paragraph{Eq.\ [7] --- Regular-Module Hilbert Space (eq:HF-regular, line 94)}
\[
  \mathcal{H}_F := A = M_d(\mathbb{C})
\]
\begin{itemize}
  \item \textbf{Computes:} The finite Hilbert space for Branch A.
  \item \textbf{Input:} $d = 64$.
  \item \textbf{Output:} $\dim(\mathcal{H}_F) = d^2 = 4096$.
  \item \textbf{Analytic/Numerical:} Analytic (definition).
\end{itemize}

\paragraph{Eq.\ [8] --- Gauge Action (line 98)}
\[
  \mathrm{ad}_a(X) = [a, X]
\]
\begin{itemize}
  \item \textbf{Computes:} Inner derivation gauge action.
  \item \textbf{Analytic/Numerical:} Analytic.
\end{itemize}

\paragraph{Eq.\ [9] --- Democratic Trace (eq:tr-dem, line 104)}
\[
  \mathrm{tr}_{\rm dem}(\cdot) := \frac{1}{d^2}\,\mathrm{Tr}_{\mathcal{H}_F}(\cdot)
\]
\begin{itemize}
  \item \textbf{Computes:} UV-normalized trace per matrix degree of freedom.
  \item \textbf{Input:} $d = 64$.
  \item \textbf{Output:} Normalization factor $1/4096$.
  \item \textbf{Analytic/Numerical:} Analytic.
\end{itemize}

\paragraph{Eq.\ [10] --- Physics (per-generator) Trace (line 111)}
\[
  \mathrm{tr}_{\mathfrak{su}}(\cdot) = \frac{1}{d^2 - 1}\,
  \mathrm{Tr}_{\mathfrak{su}}(\cdot)
\]
\begin{itemize}
  \item \textbf{Computes:} Canonical generator normalization on $\mathfrak{su}(d)$.
  \item \textbf{Input:} $d = 64$.
  \item \textbf{Output:} Normalization factor $1/4095$.
  \item \textbf{Analytic/Numerical:} Analytic.
\end{itemize}

\paragraph{Eq.\ [11] --- BOOST Factor, Branch A (eq:boost-factor, line 114) \textbf{[BOXED]}}
\[
  \boxed{\varepsilon_{\rm adj}^{(A)} = \frac{d^2}{d^2 - 1}}
\]
\begin{itemize}
  \item \textbf{Computes:} Ratio of democratic UV trace to per-generator physics trace.
        This is the ``forced binary fork'' boost factor.
  \item \textbf{Input:} $d = 64$.
  \item \textbf{Output:} $\varepsilon_{\rm adj}^{(A)} = 4096/4095 = 1.000244\ldots$
  \item \textbf{Analytic/Numerical:} Analytic (exact rational).
  \item \textbf{HOW IT ENTERS:} This multiplies the baseline $\alpha^{-1}$ to produce
        the final sub-ppm result. This is the ENTIRE content of the ``forced binary fork.''
\end{itemize}

\paragraph{Eq.\ [12] --- BOOST Factor Evaluated (line 119)}
\[
  \varepsilon_{\rm adj}^{(A)} = \frac{4096}{4095} = 1.000244\ldots
\]
\begin{itemize}
  \item \textbf{Value:} $4096/4095$.
  \item \textbf{Analytic/Numerical:} Analytic (arithmetic).
\end{itemize}

\paragraph{Eq.\ [13] --- DROP Factor, Branch B (eq:drop-factor, line 126)}
\[
  \varepsilon_{\rm adj}^{(B)} = \frac{d^2 - 1}{d^2} = \frac{4095}{4096} = 0.999756\ldots
\]
\begin{itemize}
  \item \textbf{Computes:} The alternative (falsified) trace conversion.
  \item \textbf{Value:} $4095/4096$ (reciprocal of BOOST).
\end{itemize}

\paragraph{Eq.\ [14] --- Branch A Result (Table, line 145)}
\[
  \alpha^{-1}_A = 137.03599985 \quad (\text{residual: } -0.006 \text{ ppm})
\]

\paragraph{Eq.\ [15] --- Branch B Result (Table, line 146)}
\[
  \alpha^{-1}_B = 137.03014445 \quad (\text{residual: } +42.7 \text{ ppm})
\]

\paragraph{Eq.\ [16] --- Summary Result (line 253) \textbf{[BOXED]}}
\[
  \boxed{\alpha^{-1} = 137.03599985 \quad
  (\text{residual: } {-}0.006\text{ ppm})}
\]

%-----------------------------------------------------------------------------
\subsection*{B. The Derivation Chain Table (Table 2, lines 184--207)}
%-----------------------------------------------------------------------------

The table at line 184 lists the COMPLETE derivation chain.
Listed ingredients and their values:

\begin{center}
\renewcommand{\arraystretch}{1.3}
\begin{tabular}{llll}
\toprule
Component & Value & Source & Type \\
\midrule
$K_{\mathbb{C}P^2} = \mathcal{O}(-3)$ & $-3$ & Algebraic geometry & Rigorous \\
$L_{\det} = K^{-1}$ & $\mathcal{O}(3)$ & Spin$^c$ structure & Rigorous \\
$d = k + 4$ & $64$ & $\dim H^0(\mathcal{O}(k+3))$ & Rigorous \\
$(d-1)/d$ & $63/64 = 0.984375$ & Traceless projection & Derived \\
$N_{\rm species}$ & $7$ & SM SU(2) components & SM content \\
$\mathrm{Tr}(Y^2)$ & $10$ & SM hypercharges & SM content \\
$g_F$ & $8$ & Spectral triple ($J \times \gamma \times \mathbb{C}$) & Derived \\
$w = N_{\rm species}/(g_F \cdot \mathrm{Tr}(Y^2))$ & $7/80 = 0.0875$ & Hypercharge weighting & Derived \\
$\varepsilon_{\rm adj}$ & $4096/4095$ & Regular-module trace & Forced \\
\midrule
$\alpha^{-1}$ & $137.03599985$ & All above combined & $< 0.01$ ppm \\
\bottomrule
\end{tabular}
\end{center}

%-----------------------------------------------------------------------------
\subsection*{C. CRITICAL FINDING: The Missing Master Formula}
%-----------------------------------------------------------------------------

\textbf{Section 8C does NOT contain an explicit equation of the form}
\[
  \alpha^{-1} = f(\beta_{U(1)},\; w,\; \Lambda^3,\; \varepsilon_{\rm adj},\; \ldots)
\]
\textbf{that combines the table ingredients into a single closed-form expression.}

The section provides:
\begin{enumerate}
  \item A table of ingredients (Table 2)
  \item The final numerical result $137.03599985$
  \item The statement ``All above combined''
\end{enumerate}

But the actual formula that multiplies/divides these ingredients is NEVER written
as a displayed equation in this file. The reader must reconstruct it from context
across multiple sections and appendices.

%-----------------------------------------------------------------------------
\subsection*{D. Reconstructed Computational Chain}
%-----------------------------------------------------------------------------

From appendix\_K.tex and section\_alpha\_locked.tex combined, the chain is:

\paragraph{Step 1: Fix $k_{\max} = 60$ (Analytic).}
From Bridge Lemma / Spin$^c$ index on $\mathbb{C}P^2$:
\[
  k_{\max} = \chi(\mathbb{C}P^2, \mathcal{O}(9) \oplus \mathcal{O}^{\oplus 5})
  = 55 + 5 = 60.
\]
Value: $k_{\max} = 60$.

\paragraph{Step 2: Compute $\beta_{U(1)}$ (Numerical).}
Chern--Simons weighted expectation value:
\[
  \beta_{U(1)} = \langle k + 2 \rangle
  = \frac{\displaystyle\sum_{k=0}^{59} (k+2)\, w_{\rm CS}(k)}
         {\displaystyle\sum_{k=0}^{59} w_{\rm CS}(k)},
  \qquad
  w_{\rm CS}(k) = \frac{2}{k+2}\sin^2\!\frac{\pi}{k+2}.
\]
Value: $\beta_{U(1)} = 3.7969$.

\textbf{This is the ONLY numerical (non-analytic) step in the chain.}
The sum must be evaluated computationally.

\paragraph{Step 3: Derive $d = k_{\max} + 4$ (Analytic).}
\[
  d = 60 + 4 = 64.
\]

\paragraph{Step 4: Compute finite-size factor $(d-1)/d$ (Analytic).}
\[
  \frac{d-1}{d} = \frac{63}{64} = 0.984375.
\]

\paragraph{Step 5: Compute spectral cutoff $\Lambda^3$ (Analytic).}
\[
  \Lambda^3 = k \cdot \frac{k+3}{k+4} = 60 \times \frac{63}{64} = 59.0625.
\]
\textbf{NOTE:} The text also references $\Lambda^3 = 885.9375$.
Discrepancy: $885.9375 / 59.0625 = 15$.
This factor of 15 likely arises from additional geometry
(e.g., $N_{\rm gen} \times n = 3 \times 5 = 15$,
or from the ``$a = 9, n = 5, N = 3$'' parameters mentioned at line 210).
The exact formula producing 885.9375 is not given as a displayed equation.

\paragraph{Step 6: SM Inputs (Analytic).}
\[
  N_{\rm species} = 7, \quad \mathrm{Tr}(Y^2) = 10, \quad g_F = 8.
\]
These are Standard Model content:
\begin{itemize}
  \item $N_{\rm species} = 7$: SU(2) weak isospin components per generation
        (from the SM fermion multiplet structure).
  \item $\mathrm{Tr}(Y^2) = 10$: Sum of squared hypercharges over one generation
        (using $3 \times$ generations; note appendix Z quotes $10/3$ per generation,
        so $10 = 3 \times 10/3$).
  \item $g_F = 8$: Spectral triple grading factor
        ($J \times \gamma \times \mathbb{C}$ structure).
\end{itemize}

\paragraph{Step 7: Hypercharge Weighting $w$ (Analytic).}
\[
  w = \frac{N_{\rm species}}{g_F \cdot \mathrm{Tr}(Y^2)}
  = \frac{7}{8 \times 10} = \frac{7}{80} = 0.0875.
\]

\paragraph{Step 8: BOOST Factor (Analytic).}
\[
  \varepsilon_{\rm adj} = \frac{d^2}{d^2 - 1}
  = \frac{4096}{4095} = 1.000244\ldots
\]

\paragraph{Step 9: Combine to get $\alpha^{-1}$ (THE MISSING FORMULA).}

The paper never writes a single closed-form equation.
Based on the table ingredients, the reconstructed master formula must be
something of the form:
\[
  \alpha^{-1} = \mathcal{F}\!\bigl(\beta_{U(1)},\; \Lambda^3,\; w,\;
  \varepsilon_{\rm adj}\bigr) = 137.03599985.
\]

\textbf{The explicit functional form $\mathcal{F}$ is not stated in
section\_alpha\_locked.tex.} It is presumably derived from the spectral action
expansion (the $a_4$ Seeley--DeWitt coefficient), where the gauge kinetic term
emerges as:
\[
  S_{\rm gauge} \propto f_0 \cdot a_4 \supset
  \frac{1}{4} \kappa_{U(1)} F_{\mu\nu}F^{\mu\nu},
\]
and $\alpha^{-1} = 4\pi \kappa_{U(1)}$ with $\kappa_{U(1)}$ depending on the
table ingredients.

%-----------------------------------------------------------------------------
\subsection*{E. Specific Answers to Key Questions}
%-----------------------------------------------------------------------------

\paragraph{Q1: Does the paper give $\kappa_{U(1)} = 7.26999$ anywhere?}

\textbf{NO.} The value $\kappa_{U(1)}$ does not appear numerically anywhere
in section\_alpha\_locked.tex. The file presents only the final answer
$\alpha^{-1} = 137.03599985$ and the table of ingredients.
There is no equation ``$\kappa_{U(1)} = \ldots$'' with an explicit formula
in this section.

\paragraph{Q2: Between $\beta = 3.7969$ and the final result 137.03599985,
what INTERMEDIATE values appear?}

\textbf{None.} The section jumps directly from the table of inputs to the
final result. The only intermediate values given are:
\begin{itemize}
  \item $d = 64$ (from $k + 4$)
  \item $(d-1)/d = 63/64$
  \item $w = 7/80$
  \item $\varepsilon_{\rm adj} = 4096/4095$
  \item $\Lambda^3 = 885.9375$ (stated in text, not as equation)
\end{itemize}
No intermediate coupling constant, no $\kappa_{U(1)}$, no partial product
is displayed between these table entries and the output 137.03599985.

\paragraph{Q3: Exactly HOW does the BOOST factor $4096/4095$ enter?}

The BOOST factor $\varepsilon_{\rm adj}^{(A)} = d^2/(d^2-1) = 4096/4095$
arises from the TRACE CONVERSION between:
\begin{itemize}
  \item The ``democratic'' UV trace: $\mathrm{tr}_{\rm dem} = (1/d^2)\,\mathrm{Tr}$
        (normalization per matrix degree of freedom in $M_d(\mathbb{C})$)
  \item The ``physics'' per-generator trace: $\mathrm{tr}_{\mathfrak{su}} =
        (1/(d^2-1))\,\mathrm{Tr}_{\mathfrak{su}}$ (canonical normalization
        on the Lie algebra $\mathfrak{su}(d)$)
\end{itemize}
It enters as a multiplicative correction: $\alpha^{-1}_{\rm final} =
\alpha^{-1}_{\rm baseline} \times \varepsilon_{\rm adj}$.

The ``forced binary fork'' is:
\begin{itemize}
  \item Branch A: $\mathcal{H}_F = M_d(\mathbb{C})$ (algebra itself, regular module)
        $\Rightarrow$ BOOST = $4096/4095$ $\Rightarrow$ $\alpha^{-1} = 137.03599985$
  \item Branch B: $\mathcal{H}_F = \mathbb{C}^d$ (fermion representation)
        $\Rightarrow$ DROP = $4095/4096$ $\Rightarrow$ $\alpha^{-1} = 137.03014445$
\end{itemize}
The fork is ``forced'' because once you choose the Hilbert space, the trace
normalization is determined with no remaining freedom. Branch B misses by
43 ppm and cannot be rescued.

\paragraph{Q4: Is there a line ``$\kappa_{U(1)} = \ldots$'' with an explicit formula?}

\textbf{NO.} Section 8C does not contain any equation of the form
$\kappa_{U(1)} = (\text{something})$. The coupling $\kappa_{U(1)}$ is
mentioned indirectly via the stiffness ratio
$\kappa_{U(1)}/\kappa_{SU(2)} = n_1/n_2 = 1/2$
(in appendix\_K.tex, line 127), but its absolute value is never stated.

\paragraph{Q5: Is the step from $\beta_{U(1)} = 3.7969$ to
$\alpha^{-1} = 137.036$ analytic or numerical?}

\textbf{Hybrid.} The computation of $\beta_{U(1)}$ from the CS sum is
numerical (60-term weighted sum). All subsequent steps (finite-size factor,
SM content, BOOST) are analytic/arithmetic. But the formula connecting
$\beta_{U(1)}$ to $\alpha^{-1}$ is not displayed.

%-----------------------------------------------------------------------------
\subsection*{F. What Is Actually in section\_alpha\_locked.tex vs.\ What Is Missing}
%-----------------------------------------------------------------------------

\paragraph{Present:}
\begin{enumerate}
  \item Operator choice (connection Laplacian, locked)
  \item Regularization rule (plateau cutoff, locked)
  \item Finite-$k$ truncation giving $d = k + 4 = 64$
  \item Spectral cutoff $\Lambda^3 = k(k+3)/(k+4)$
  \item The forced binary fork (regular-module vs.\ fermion-rep)
  \item BOOST factor $= 4096/4095$ (derived, exact)
  \item DROP factor $= 4095/4096$ (the rejected alternative)
  \item Table of all ingredients with values
  \item Final numerical result $\alpha^{-1} = 137.03599985$
  \item The SM inputs $N_{\rm species} = 7$, $\mathrm{Tr}(Y^2) = 10$, $g_F = 8$
  \item Hypercharge weighting $w = 7/80$
  \item Falsification criteria
\end{enumerate}

\paragraph{Missing:}
\begin{enumerate}
  \item The EXPLICIT master formula combining ingredients into $\alpha^{-1}$
  \item The absolute value of $\kappa_{U(1)}$
  \item Any intermediate numerical products
        (e.g., ``$\kappa_{U(1)} = \beta \times \Lambda^3 \times w \times
        \varepsilon / (4\pi) = \ldots$'')
  \item The derivation of $\Lambda^3 = 885.9375$ from the stated parameters
        $a = 9, n = 5, N = 3$ (only stated as text, no formula)
  \item The connection between the lattice/CS computation producing
        $\beta_{U(1)} = 3.7969$ and the perturbative/spectral-action
        framework using $\kappa_{U(1)}$
\end{enumerate}

%-----------------------------------------------------------------------------
\subsection*{G. Numerical Verification Attempt}
%-----------------------------------------------------------------------------

Working backward from the result:
\[
  \alpha^{-1} = 137.03599985
  \quad\Rightarrow\quad
  \alpha^{-1}/\varepsilon_{\rm adj} = 137.03599985 / (4096/4095)
  = 137.03599985 \times (4095/4096) = 137.00253\ldots
\]
So the ``baseline'' $\alpha^{-1}$ before the BOOST is approximately $137.003$.

If $\alpha^{-1} = 4\pi\kappa_{U(1)}$, then:
\[
  \kappa_{U(1)} = 137.03599985 / (4\pi) = 10.9057\ldots
\]
And the baseline would be $\kappa_{U(1),\rm base} = 137.003/(4\pi) = 10.903$.

These intermediates are NEVER stated in the paper.

%-----------------------------------------------------------------------------
\subsection*{H. Summary Assessment}
%-----------------------------------------------------------------------------

Section 8C (section\_alpha\_locked.tex) is a \textbf{convention-locking} document,
not a \textbf{derivation} document. Its purpose is to:
\begin{enumerate}
  \item Specify which operator, regulator, and microsector trace are used (locking
        all conventions)
  \item Present the forced binary fork and show that only Branch A survives
  \item Tabulate all ingredients and state the final result
\end{enumerate}

The actual derivation --- the formula connecting the ingredients to $\alpha^{-1}$
--- is distributed across multiple files:
\begin{itemize}
  \item appendix\_K.tex: CS level sum giving $\beta_{U(1)} = 3.7969$
  \item The spectral action / heat kernel framework (producing $\kappa$ from $a_4$)
  \item section\_alpha\_locked.tex: The BOOST correction and SM inputs
\end{itemize}

\textbf{The computational chain has a ``black box'' between the table of inputs
and the output 137.03599985.} The paper asserts the result but does not show the
intermediate multiplication in any single equation.

\end{document}
